% =====================================================================
%  DAY 1 — 2026 Comprehensive Exam
%  Ryan Gibbons
%
%  HOW TO USE
%    Write the Day 1 response in the body below.
%    Build:  latexmk -pdf "Day 1 Gibbons.tex"
%      or:   python3 tools/build_all.py "Day 1*"
%
%  CITATIONS — American Political Science Review style.
%    \cite{key}            -> (Putnam 1993)          parenthetical
%    \citeasnoun{key}      -> Putnam (1993)          narrative / in-text
%    \citeyear{key}        -> 1993                   year alone
%    \cite[45]{key}        -> (Putnam 1993, 45)     with page number
%    \possessivecite{key}  -> Putnam's (1993)        possessive
%    Multiple: \cite{a,b,c}
%    Keys come from ../refs.bib. Nothing is cited yet — that is fine;
%    the References section simply stays empty until you cite something.
% =====================================================================


% ---------------------------------------------------------------------
%  SETTING 1 of 3 — FONT SIZE
%  Change the option on the next line. Valid values: 10pt, 11pt, 12pt.
% ---------------------------------------------------------------------
\documentclass[12pt]{article}


% ---------------------------------------------------------------------
%  SETTING 2 of 3 — MARGINS
%  One value applied to all four sides. Examples: 1in, 1.25in, 2.5cm.
%  For uneven margins, replace the \geometry line further down.
% ---------------------------------------------------------------------
\newcommand{\MarginSize}{1in}

% ---------------------------------------------------------------------
%  SETTING 3 of 3 — LINE SPACING
%  1.0 = single, 1.5 = one-and-a-half, 2.0 = double. Decimals are fine.
% ---------------------------------------------------------------------
\newcommand{\LineSpacing}{2.0}


% ---------------------------------------------------------------------
%  TITLE PAGE FIELDS — edit these four lines per response
% ---------------------------------------------------------------------
\newcommand{\PaperTitle}{Day 1}
\newcommand{\PaperAuthor}{Ryan Gibbons}
\newcommand{\PaperDate}{\today}
\newcommand{\PaperSubtitle}{2026 Comparative Politics Comprehensive Exam \\[0.3em] Day 1: Core}


% =====================================================================
%  PACKAGES
% =====================================================================

% --- Page layout, spacing, headers ---
\usepackage[margin=\MarginSize]{geometry}
\usepackage{setspace}
\usepackage{fancyhdr}
\usepackage{titlesec}
\usepackage{titling}

% --- Fonts and text ---
% lmodern is Computer Modern redrawn as a scalable outline font: it looks the
% same as the LaTeX default but, unlike the default, microtype can expand it
% and PDF text search/copy works properly.
\usepackage{lmodern}
\usepackage[T1]{fontenc}
\usepackage[utf8]{inputenc}
\usepackage{microtype}      % subtle spacing improvements
\usepackage{csquotes}       % \enquote{...} for nested quotation marks
\usepackage{soul}           % \hl{} highlight, \ul{} underline
\usepackage{xcolor}
\usepackage{url}

% --- Mathematics ---
\usepackage{amsmath}
\usepackage{amssymb}
\usepackage{amsfonts}
\usepackage{amsthm}
\usepackage{mathtools}      % \coloneqq, \DeclarePairedDelimiter, etc.
\usepackage{bm}             % \bm{} bold math
\usepackage{dsfont}         % \mathds{1} indicator function

% --- Figures ---
\usepackage{graphicx}
\usepackage{float}          % [H] to pin a float exactly here
\usepackage{subcaption}
\usepackage{caption}
\usepackage{wrapfig}
\usepackage{rotating}

% --- Tables ---
\usepackage{booktabs}       % \toprule \midrule \bottomrule
\usepackage{tabularx}
\usepackage{longtable}
\usepackage{multirow}
\usepackage{array}
\usepackage{threeparttable}
\usepackage{siunitx}        % \num{1234.5}, aligned numeric columns

% --- Lists, code, diagrams ---
\usepackage{enumitem}
\usepackage{listings}
\usepackage{tikz}
\usetikzlibrary{shapes.geometric, arrows, arrows.meta, calc, positioning, trees}
\usepackage{pgfplots}
\pgfplotsset{compat=1.18}

% --- Citations: APSR style (harvard.sty + apsr.bst, both in this folder) ---
\usepackage[apsr]{harvard}

% --- Cross-references. Loaded last, as hyperref prefers to be. ---
\usepackage{hyperref}
\hypersetup{
  colorlinks = true,
  linkcolor  = black,      % internal links: black for a printed draft
  citecolor  = black,
  urlcolor   = blue,
  pdftitle   = {\PaperTitle},
  pdfauthor  = {\PaperAuthor},
}


% =====================================================================
%  APPLY SETTINGS
% =====================================================================
\setstretch{\LineSpacing}

% Page numbers bottom-right, no rules, no header text.
\pagestyle{fancy}
\fancyhf{}
\rfoot{\thepage}
\renewcommand{\headrulewidth}{0pt}
\renewcommand{\footrulewidth}{0pt}

% Paragraphs: indented, no extra vertical gap (standard for prose).
% For block paragraphs instead, set parindent to 0pt and parskip to ~1em.
\setlength{\parindent}{0.5in}
\setlength{\parskip}{0pt}


% =====================================================================
%  THEOREM ENVIRONMENTS
% =====================================================================
\theoremstyle{plain}
\newtheorem{theorem}{Theorem}
\newtheorem{prop}{Proposition}
\newtheorem{lemma}{Lemma}
\newtheorem{corollary}{Corollary}

\theoremstyle{definition}
\newtheorem{definition}{Definition}
\newtheorem{assumption}{Assumption}
\newtheorem{example}{Example}

\theoremstyle{remark}
\newtheorem*{remark}{Remark}


% =====================================================================
%  DOCUMENT
% =====================================================================
\begin{document}

% --- Title page (no page number) ---
\begin{titlepage}
    \centering
    \vspace*{2cm}
    {\large \PaperSubtitle \par}
    \vspace{8cm}
    {\large \PaperAuthor \par}
    \vspace{.2cm}
    {\large University of Chicago \par}
    \vspace{.2cm}
    {\large Department of Political Science \par}
    \vspace{1cm}
    {\large August 22, 2026 \par}
    \vfill
\end{titlepage}

% Body starts here; page numbering restarts at 1.
\newpage

\vspace{2cm}

\section*{Question}

What array of causal drivers are adduced in the literature to account for different political outcomes in regime type, state form or electoral systems (i.e. dictatorship vs democracy; contrasting cleavage structures; electoral systems). How do theorists deploy history and timing in making their arguments?

\thispagestyle{empty}


\newpage
\setcounter{page}{1}

\section*{Introduction}


From the peninsula of Gibraltar,\footnote{So long as one's southward view is not entirely obscured by the titular Rock} one can see land governed by three distinct states: the Kingdom of Spain, the Kingdom of Morocco, and the United Kingdom.\footnote{Of which Gribraltar is an ``Overseas Territory''} The landscape is visually quite similar; the boundary between Gibraltar and Spain is purely political, and the seven-mile-wide strait is geologically little more than an advanced valley. Yet, each of these lands has experienced dramatically distinct transitions -- or lack thereof -- from monarchy to democracy. Spain's development to liberal democracy famously included an extensive fascist dictatorship; the United Kingdom's gradual progression to the current Westminster System was significantly more gradual; Morocco is still a monarchy today, even if with some modest constitutional restrictions. Why, then, does this one coastal viewpoint offer such a distinct range of political outcomes across a relatively undestinctive physical landscape? \\

Identifying political and economic deviations across meager physical boundaries is not a novel perspective within political science or the social sciences more generally; Açemoglu \& Robinson \citeyear{acemoglurobinson2012} famously employ this framework to counter theories of economic dependence on geographical factors in ``Why Nations Fail.'' Yet, the question is still an intriguing one that has motivated a significant corner of the social sciences, perhaps none more significantly than comparative politics. By Mill's \citeyear{mill1843} ``method of difference,'' the relatively identical nature of this physical landscape, compared to the significant variation in sociopolitical outcomes, suggests that the causal factors behind regime outcomes must lie elsewhere. If geography alone does not sufficiently explain the distinction between the human-created institutions built upon them, then it indeed must be the case that the driving factors are, to a large extent, humanistic themselves. \\

In this response, I explore the literature's fundamental discussions of the factors driving political outcomes in regime type, state form, and broadly, political institutions. I present three primary causal arguments, each of which is unique in its approach but not mutually exclusive from the others: (1) a state's historical economic composition and structure; (2) the identity-based composition of the constituency, along with the salience of those identity characteristics; and (3) the chronological timing of modernization processes themselves. I further present a relatively -- though admittedly not entirely -- conflicting counterargument regarding the innate dependency of political outcomes on randomized or exogenous influences, discussing both its limitations and its possible contributions nonetheless. I close by exploring some various empirical strategies that could be used to further our collective perception of these three (or four) causal relationships. \textcolor{red}{rewrite this -- explore existing and present further} \\

In discussing and evaluating these arguments, I do not intend to provide a definitive weighting of their relative accuracy or value. To do so would be more than I could hope to reasonably accomplish in an entire career, and would almost certainly involve discarding valuable arguments that may indeed be foundational. Yet, I do explore my own perspectives within these frameworks, largely built upon my personal affinity towards rationalist and institutionalist answers. It is our responsibility as social scientists to pursue those narratives which we believe most effectively explain the world around us, and to consider how those mechanisms might be best employed to construct a more effective political future.


\section*{An Economic Explanation} 








% =====================================================================
%  REFERENCES
%  Draws on two bibliographies at the repo root:
%    ../refs.bib    hand-curated. Edit freely; never overwritten.
%    ../zotero.bib  the whole Zotero library, generated by
%                   tools/export_zotero_bib.py. Do not edit by hand.
%  Look a key up in ../zotero-keys.txt, or just search ../zotero.bib.
%  Uncomment \nocite{*} to print every entry in refs.bib regardless of
%  whether it is cited — useful for checking the file, not for a draft.
% =====================================================================
\newpage
% \nocite{*}

% UNCOMMENT OUT ONCE STARTED
\bibliography{../refs,../zotero}

\end{document}
